% --- Archon physics formalization source begin ---
% archon:physics
% archon:covers PhyXMiniProblems/problem_phyx_mini_0985.lean
% archon:source-report reports/phyx_mini/problem_phyx_mini_0985.source.json

\chapter{Physics problem phyx\_mini\_0985}
\label{ch:PhyXMiniProblems_problem_phyx_mini_0985}

\paragraph{Problem source.}
\#\# Physical scenario

An inductor with inductance \textbackslash{}( L = 0.200\textbackslash{} \textbackslash{}text\{H\} \textbackslash{}) and negligible resistance is connected to a battery, a switch \textbackslash{}( S \textbackslash{}), and two resistors, \textbackslash{}( R\_1 = 8.00\textbackslash{} \textbackslash{}Omega \textbackslash{}) and \textbackslash{}( R\_2 = 6.00\textbackslash{} \textbackslash{}Omega \textbackslash{}). The battery has emf \textbackslash{}( 48.0\textbackslash{} \textbackslash{}text\{V\} \textbackslash{}) and negligible internal resistance. \textbackslash{}( S \textbackslash{}) is closed at \textbackslash{}( t = 0 \textbackslash{}).

\#\# Question

When \textbackslash{}( i\_3 \textbackslash{}) has half of its final value, what is \textbackslash{}( i\_1 \textbackslash{})?s

\#\# Answer choices

- A: A: 3.00A

- B: B: 4.71A

- C: C: 6.00A

- D: D: 4.00A

\#\# Figure caption (auxiliary; use the image as primary evidence)

The image depicts an electrical circuit diagram consisting of the following components:

1. **Battery (\textbackslash{}( \textbackslash{}mathcal\{E\} \textbackslash{}))**: Located on the left, it provides the electromotive force (\textbackslash{}(+\textbackslash{})) and drives the current \textbackslash{}(i\_1\textbackslash{}).

2. **Switch (\textbackslash{}(S\textbackslash{}))**: Positioned right after the battery in the circuit, which can open or close the circuit.

3. **Resistor \textbackslash{}(R\_1\textbackslash{})**: Connected in series with the switch, through which current \textbackslash{}(i\_2\textbackslash{}) flows.

4. **Resistor \textbackslash{}(R\_2\textbackslash{})**: Arranged in parallel with \textbackslash{}(R\_1\textbackslash{}). The current \textbackslash{}(i\_3\textbackslash{}) flows through this resistor.

5. **Inductor (\textbackslash{}(L\textbackslash{}))**: Connected in series with \textbackslash{}(R\_2\textbackslash{}).

6. **Current Arrows**:
   - \textbackslash{}(i\_1\textbackslash{}): Enters the circuit from the battery, shown with an upward arrow.
   - \textbackslash{}(i\_2\textbackslash{}): Flows through \textbackslash{}(R\_1\textbackslash{}), directed downward.
   - \textbackslash{}(i\_3\textbackslash{}): Flows through the inductor \textbackslash{}(L\textbackslash{}), also directed downward.

The layout suggests a basic RLC circuit, with the relationships between the components highlighting series and parallel arrangements and current flow directions.

\#\# Dataset metadata

Electric Circuits; Multi-Formula Reasoning, Numerical Reasoning

\paragraph{Recorded answer/context.}
B: B: 4.71A

\paragraph{Figure/image path.}
/root/proposal\_for\_physic/science-mango/phyx\_mini\_run/phyx\_data/test\_image/985.png

\paragraph{Formalization target.}
create a compiling Lean file with sorry bodies at `PhyXMiniProblems/problem\_phyx\_mini\_0985.lean`. The Lean declarations must preserve the physical quantities, dimensions or dimensional roles, figure labels, governing-law hypotheses, and final relation expressed by this problem.
Use Mathlib/Physlib names found through LeanExplore where available. If a physics API is missing, introduce faithful local abstractions rather than scalar placeholder aliases.

\begin{theorem}[Physics formalization target]
\label{thm:physics:phyx_mini_0985:target}
\lean{PhyXMiniProblems.ProblemPhyXMini0985.problem_phyx_mini_0985}
\uses{def:physics:phyx-mini-0985:phyxminiproblems-problemphyxmini0985-electriccurrentinamperes, def:physics:phyx-mini-0985:phyxminiproblems-problemphyxmini0985-parallelrltransientsetup, def:physics:phyx-mini-0985:phyxminiproblems-problemphyxmini0985-matchesidealparallelrlscenario, def:physics:phyx-mini-0985:phyxminiproblems-problemphyxmini0985-matchessuppliedparallelrlfigure, def:physics:phyx-mini-0985:phyxminiproblems-problemphyxmini0985-matchesgivenparallelrldata, def:physics:phyx-mini-0985:phyxminiproblems-problemphyxmini0985-satisfiesidealparallelrltransientlaws, def:physics:phyx-mini-0985:phyxminiproblems-problemphyxmini0985-recordeddatasetanswer}
The assigned autoformalize agent should translate this physics problem into Lean declarations in the covered file, with theorem and lemma proof bodies written as `by sorry`.
\end{theorem}
\begin{proof}
This is an autoformalization task, not a proof task. Produce statements that can later be proved by the physics prover without weakening the physical model.
\end{proof}

% --- Archon declaration topology begin ---
\section{Declaration topology}
\begin{definition}[electric Current Dimension]
  \label{def:physics:phyx-mini-0985:phyxminiproblems-problemphyxmini0985-electriccurrentdimension}
  \lean{PhyXMiniProblems.ProblemPhyXMini0985.electricCurrentDimension}
  Electric current has physical dimension charge per time
\end{definition}
\begin{proof}
  This is the declared model definition of electric Current Dimension.
\end{proof}
\begin{definition}[electric Potential Dimension]
  \label{def:physics:phyx-mini-0985:phyxminiproblems-problemphyxmini0985-electricpotentialdimension}
  \lean{PhyXMiniProblems.ProblemPhyXMini0985.electricPotentialDimension}
  Electric potential difference has physical dimension energy per charge
\end{definition}
\begin{proof}
  This is the declared model definition of electric Potential Dimension.
\end{proof}
\begin{definition}[electrical Resistance Dimension]
  \label{def:physics:phyx-mini-0985:phyxminiproblems-problemphyxmini0985-electricalresistancedimension}
  \lean{PhyXMiniProblems.ProblemPhyXMini0985.electricalResistanceDimension}
  \uses{def:physics:phyx-mini-0985:phyxminiproblems-problemphyxmini0985-electriccurrentdimension, def:physics:phyx-mini-0985:phyxminiproblems-problemphyxmini0985-electricpotentialdimension}
  Electrical resistance has dimension potential difference per current
\end{definition}
\begin{proof}
  This is the declared model definition of electrical Resistance Dimension.
\end{proof}
\begin{definition}[electrical Inductance Dimension]
  \label{def:physics:phyx-mini-0985:phyxminiproblems-problemphyxmini0985-electricalinductancedimension}
  \lean{PhyXMiniProblems.ProblemPhyXMini0985.electricalInductanceDimension}
  \uses{def:physics:phyx-mini-0985:phyxminiproblems-problemphyxmini0985-electricalresistancedimension}
  Inductance has dimension resistance times time
\end{definition}
\begin{proof}
  This is the declared model definition of electrical Inductance Dimension.
\end{proof}
\begin{definition}[Electric Current]
  \label{def:physics:phyx-mini-0985:phyxminiproblems-problemphyxmini0985-electriccurrent}
  \lean{PhyXMiniProblems.ProblemPhyXMini0985.ElectricCurrent}
  \uses{def:physics:phyx-mini-0985:phyxminiproblems-problemphyxmini0985-electriccurrentdimension}
  A signed, unit-independent branch current
\end{definition}
\begin{proof}
  This is the declared model definition of Electric Current.
\end{proof}
\begin{definition}[Electric Potential Difference]
  \label{def:physics:phyx-mini-0985:phyxminiproblems-problemphyxmini0985-electricpotentialdifference}
  \lean{PhyXMiniProblems.ProblemPhyXMini0985.ElectricPotentialDifference}
  \uses{def:physics:phyx-mini-0985:phyxminiproblems-problemphyxmini0985-electricpotentialdimension}
  A signed, oriented, unit-independent potential difference
\end{definition}
\begin{proof}
  This is the declared model definition of Electric Potential Difference.
\end{proof}
\begin{definition}[Electrical Resistance Magnitude]
  \label{def:physics:phyx-mini-0985:phyxminiproblems-problemphyxmini0985-electricalresistancemagnitude}
  \lean{PhyXMiniProblems.ProblemPhyXMini0985.ElectricalResistanceMagnitude}
  \uses{def:physics:phyx-mini-0985:phyxminiproblems-problemphyxmini0985-electricalresistancedimension}
  A nonnegative, unit-independent resistance magnitude
\end{definition}
\begin{proof}
  This is the declared model definition of Electrical Resistance Magnitude.
\end{proof}
\begin{definition}[Electrical Inductance Magnitude]
  \label{def:physics:phyx-mini-0985:phyxminiproblems-problemphyxmini0985-electricalinductancemagnitude}
  \lean{PhyXMiniProblems.ProblemPhyXMini0985.ElectricalInductanceMagnitude}
  \uses{def:physics:phyx-mini-0985:phyxminiproblems-problemphyxmini0985-electricalinductancedimension}
  A nonnegative, unit-independent inductance magnitude
\end{definition}
\begin{proof}
  This is the declared model definition of Electrical Inductance Magnitude.
\end{proof}
\begin{definition}[electric Current In Amperes]
  \label{def:physics:phyx-mini-0985:phyxminiproblems-problemphyxmini0985-electriccurrentinamperes}
  \lean{PhyXMiniProblems.ProblemPhyXMini0985.electricCurrentInAmperes}
  \uses{def:physics:phyx-mini-0985:phyxminiproblems-problemphyxmini0985-electriccurrent}
  Read a signed electric current in coherent SI amperes
\end{definition}
\begin{proof}
  This is the declared model definition of electric Current In Amperes.
\end{proof}
\begin{definition}[potential Difference In Volts]
  \label{def:physics:phyx-mini-0985:phyxminiproblems-problemphyxmini0985-potentialdifferenceinvolts}
  \lean{PhyXMiniProblems.ProblemPhyXMini0985.potentialDifferenceInVolts}
  \uses{def:physics:phyx-mini-0985:phyxminiproblems-problemphyxmini0985-electricpotentialdifference}
  Read an oriented potential difference or emf in coherent SI volts
\end{definition}
\begin{proof}
  This is the declared model definition of potential Difference In Volts.
\end{proof}
\begin{definition}[resistance In Ohms]
  \label{def:physics:phyx-mini-0985:phyxminiproblems-problemphyxmini0985-resistanceinohms}
  \lean{PhyXMiniProblems.ProblemPhyXMini0985.resistanceInOhms}
  \uses{def:physics:phyx-mini-0985:phyxminiproblems-problemphyxmini0985-electricalresistancemagnitude}
  Read a resistance magnitude in coherent SI ohms
\end{definition}
\begin{proof}
  This is the declared model definition of resistance In Ohms.
\end{proof}
\begin{definition}[inductance In Henries]
  \label{def:physics:phyx-mini-0985:phyxminiproblems-problemphyxmini0985-inductanceinhenries}
  \lean{PhyXMiniProblems.ProblemPhyXMini0985.inductanceInHenries}
  \uses{def:physics:phyx-mini-0985:phyxminiproblems-problemphyxmini0985-electricalinductancemagnitude}
  Read an inductance magnitude in coherent SI henries
\end{definition}
\begin{proof}
  This is the declared model definition of inductance In Henries.
\end{proof}
\begin{definition}[Circuit Node]
  \label{def:physics:phyx-mini-0985:phyxminiproblems-problemphyxmini0985-circuitnode}
  \lean{PhyXMiniProblems.ProblemPhyXMini0985.CircuitNode}
  Electrically distinct nodes needed to describe the raster's topology
\end{definition}
\begin{proof}
  This is the declared model definition of Circuit Node.
\end{proof}
\begin{definition}[Circuit Element]
  \label{def:physics:phyx-mini-0985:phyxminiproblems-problemphyxmini0985-circuitelement}
  \lean{PhyXMiniProblems.ProblemPhyXMini0985.CircuitElement}
  The five lumped elements shown in the circuit
\end{definition}
\begin{proof}
  This is the declared model definition of Circuit Element.
\end{proof}
\begin{definition}[Current Label]
  \label{def:physics:phyx-mini-0985:phyxminiproblems-problemphyxmini0985-currentlabel}
  \lean{PhyXMiniProblems.ProblemPhyXMini0985.CurrentLabel}
  Current labels printed beside the three arrows in the raster
\end{definition}
\begin{proof}
  This is the declared model definition of Current Label.
\end{proof}
\begin{definition}[Switch State]
  \label{def:physics:phyx-mini-0985:phyxminiproblems-problemphyxmini0985-switchstate}
  \lean{PhyXMiniProblems.ProblemPhyXMini0985.SwitchState}
  Idealized state of the switching element
\end{definition}
\begin{proof}
  This is the declared model definition of Switch State.
\end{proof}
\begin{definition}[Component Model]
  \label{def:physics:phyx-mini-0985:phyxminiproblems-problemphyxmini0985-componentmodel}
  \lean{PhyXMiniProblems.ProblemPhyXMini0985.ComponentModel}
  Constitutive role assigned to each displayed circuit symbol
\end{definition}
\begin{proof}
  This is the declared model definition of Component Model.
\end{proof}
\begin{definition}[Parallel RLFigure]
  \label{def:physics:phyx-mini-0985:phyxminiproblems-problemphyxmini0985-parallelrlfigure}
  \lean{PhyXMiniProblems.ProblemPhyXMini0985.ParallelRLFigure}
  \uses{def:physics:phyx-mini-0985:phyxminiproblems-problemphyxmini0985-circuitnode, def:physics:phyx-mini-0985:phyxminiproblems-problemphyxmini0985-circuitelement, def:physics:phyx-mini-0985:phyxminiproblems-problemphyxmini0985-currentlabel, def:physics:phyx-mini-0985:phyxminiproblems-problemphyxmini0985-switchstate}
  Figure-derived labels, incidences, and directed branch paths
\end{definition}
\begin{proof}
  This is the declared model definition of Parallel RLFigure.
\end{proof}
\begin{definition}[Parallel RLTransient Setup]
  \label{def:physics:phyx-mini-0985:phyxminiproblems-problemphyxmini0985-parallelrltransientsetup}
  \lean{PhyXMiniProblems.ProblemPhyXMini0985.ParallelRLTransientSetup}
  \uses{def:physics:phyx-mini-0985:phyxminiproblems-problemphyxmini0985-electriccurrent, def:physics:phyx-mini-0985:phyxminiproblems-problemphyxmini0985-electricpotentialdifference, def:physics:phyx-mini-0985:phyxminiproblems-problemphyxmini0985-electricalresistancemagnitude, def:physics:phyx-mini-0985:phyxminiproblems-problemphyxmini0985-electricalinductancemagnitude, def:physics:phyx-mini-0985:phyxminiproblems-problemphyxmini0985-circuitelement, def:physics:phyx-mini-0985:phyxminiproblems-problemphyxmini0985-currentlabel, def:physics:phyx-mini-0985:phyxminiproblems-problemphyxmini0985-switchstate, def:physics:phyx-mini-0985:phyxminiproblems-problemphyxmini0985-componentmodel, def:physics:phyx-mini-0985:phyxminiproblems-problemphyxmini0985-parallelrlfigure}
  Independent physical quantities and transient observables for the switched circuit. In particular, every current and final current is an independent field constrained below by physical laws; no target current is defined from the desired answer
\end{definition}
\begin{proof}
  This is the declared model definition of Parallel RLTransient Setup.
\end{proof}
\begin{definition}[Matches Ideal Parallel RLScenario]
  \label{def:physics:phyx-mini-0985:phyxminiproblems-problemphyxmini0985-matchesidealparallelrlscenario}
  \lean{PhyXMiniProblems.ProblemPhyXMini0985.MatchesIdealParallelRLScenario}
  \uses{def:physics:phyx-mini-0985:phyxminiproblems-problemphyxmini0985-parallelrltransientsetup}
  The standard ideal lumped-component interpretation of the five symbols
\end{definition}
\begin{proof}
  This is the declared model definition of Matches Ideal Parallel RLScenario.
\end{proof}
\begin{definition}[Matches Supplied Parallel RLFigure]
  \label{def:physics:phyx-mini-0985:phyxminiproblems-problemphyxmini0985-matchessuppliedparallelrlfigure}
  \lean{PhyXMiniProblems.ProblemPhyXMini0985.MatchesSuppliedParallelRLFigure}
  \uses{def:physics:phyx-mini-0985:phyxminiproblems-problemphyxmini0985-parallelrltransientsetup}
  Topology and arrow directions read from the primary image. The source image shows the switch symbol open even though the problem asks for its subsequent closed-circuit transient
\end{definition}
\begin{proof}
  This is the declared model definition of Matches Supplied Parallel RLFigure.
\end{proof}
\begin{definition}[Matches Given Parallel RLData]
  \label{def:physics:phyx-mini-0985:phyxminiproblems-problemphyxmini0985-matchesgivenparallelrldata}
  \lean{PhyXMiniProblems.ProblemPhyXMini0985.MatchesGivenParallelRLData}
  \uses{def:physics:phyx-mini-0985:phyxminiproblems-problemphyxmini0985-electriccurrentinamperes, def:physics:phyx-mini-0985:phyxminiproblems-problemphyxmini0985-potentialdifferenceinvolts, def:physics:phyx-mini-0985:phyxminiproblems-problemphyxmini0985-resistanceinohms, def:physics:phyx-mini-0985:phyxminiproblems-problemphyxmini0985-inductanceinhenries, def:physics:phyx-mini-0985:phyxminiproblems-problemphyxmini0985-parallelrltransientsetup}
  Numerical component data and switching protocol stated in the problem
\end{definition}
\begin{proof}
  This is the declared model definition of Matches Given Parallel RLData.
\end{proof}
\begin{definition}[Satisfies Ideal Parallel RLTransient Laws]
  \label{def:physics:phyx-mini-0985:phyxminiproblems-problemphyxmini0985-satisfiesidealparallelrltransientlaws}
  \lean{PhyXMiniProblems.ProblemPhyXMini0985.SatisfiesIdealParallelRLTransientLaws}
  \uses{def:physics:phyx-mini-0985:phyxminiproblems-problemphyxmini0985-electriccurrentinamperes, def:physics:phyx-mini-0985:phyxminiproblems-problemphyxmini0985-potentialdifferenceinvolts, def:physics:phyx-mini-0985:phyxminiproblems-problemphyxmini0985-resistanceinohms, def:physics:phyx-mini-0985:phyxminiproblems-problemphyxmini0985-inductanceinhenries, def:physics:phyx-mini-0985:phyxminiproblems-problemphyxmini0985-parallelrltransientsetup}
  Ideal transient and DC-final-value laws. These are general relations among independent observables. They do not assign the requested value to i₁
\end{definition}
\begin{proof}
  This is the declared model definition of Satisfies Ideal Parallel RLTransient Laws.
\end{proof}
\begin{definition}[Answer Choice]
  \label{def:physics:phyx-mini-0985:phyxminiproblems-problemphyxmini0985-answerchoice}
  \lean{PhyXMiniProblems.ProblemPhyXMini0985.AnswerChoice}
  Labels of the four current-valued choices printed with the problem
\end{definition}
\begin{proof}
  This is the declared model definition of Answer Choice.
\end{proof}
\begin{definition}[displayed Current In Amperes]
  \label{def:physics:phyx-mini-0985:phyxminiproblems-problemphyxmini0985-answerchoice-displayedcurrentinamperes}
  \lean{PhyXMiniProblems.ProblemPhyXMini0985.AnswerChoice.displayedCurrentInAmperes}
  \uses{def:physics:phyx-mini-0985:phyxminiproblems-problemphyxmini0985-answerchoice}
  Current printed beside each answer choice, in amperes
\end{definition}
\begin{proof}
  This is the declared model definition of displayed Current In Amperes.
\end{proof}
\begin{definition}[matches Rounded Current]
  \label{def:physics:phyx-mini-0985:phyxminiproblems-problemphyxmini0985-answerchoice-matchesroundedcurrent}
  \lean{PhyXMiniProblems.ProblemPhyXMini0985.AnswerChoice.matchesRoundedCurrent}
  \uses{def:physics:phyx-mini-0985:phyxminiproblems-problemphyxmini0985-answerchoice}
  Agreement with a displayed current rounded to the nearest hundredth
\end{definition}
\begin{proof}
  This is the declared model definition of matches Rounded Current.
\end{proof}
\begin{definition}[recorded Dataset Answer]
  \label{def:physics:phyx-mini-0985:phyxminiproblems-problemphyxmini0985-recordeddatasetanswer}
  \lean{PhyXMiniProblems.ProblemPhyXMini0985.recordedDatasetAnswer}
  \uses{def:physics:phyx-mini-0985:phyxminiproblems-problemphyxmini0985-answerchoice}
  The answer label recorded in the supplied dataset metadata
\end{definition}
\begin{proof}
  This is the declared model definition of recorded Dataset Answer.
\end{proof}
% --- Archon declaration topology end ---
% --- Archon physics formalization source end ---
